\storyheader{Cảm ơn nhé, người bạn không tồn tại của tôi
}{Ngày 20}{

Tôi ngồi ở góc lớp, chầm chậm mở từng trang sách trước mắt ra, ngồi đọc
lẩm bẩm một mình dù hiện tại đang là giờ ra chơi, khi mà tất cả mọi
người trong lớp đều đang cùng bạn bè của họ trò chuyện rôm rả.

Tôi không thật sự thích việc đọc sách cho lắm, tuy nhiên ngoài đọc sách
thì thật sự không có quá nhiều thứ để làm khi ở một mình, nhất là với
một đứa u ám và chả có một sở thích nào cụ thể như tôi.

Nếu được chọn, tôi cũng muốn được chạy ra trò chuyện và chơi đùa cùng
bọn họ.

Chỉ tiếc là, tôi không thể.

Với một người u ám như tôi, có lẽ sẽ chẳng ai muốn trò chuyện đâu.

``Hôm nay cậu ấy lại ngồi đọc sách một mình nữa nhỉ?'' Giọng của một
người bạn cùng lớp bất chợt vang lên.

Hai cậu ấy đang... nói về tôi nhỉ?

``Ừm, nghe mọi người bảo cậu ta bị bệnh tâm lý.'' Một người khác đáp lới
với một giọng điệu dè chừng, dường như đang sợ tôi nghe thấy. ``Nghe bảo
là bị điên đấy.''

``Này! Nói vậy là thô lỗ đấy!''

Đây không phải là lần đầu tôi nghe thấy mấy cuộc trò chuyện kiểu này,
với một người cứ mang một dáng vẻ u ám, gia cảnh có chút đặc biệt, khuôn
mặt ảm đảm và ít nói như tôi thì tất nhiên sẽ rộ lên mấy cái tin đồn kỳ
lạ vì bệnh tâm lý rồi.

Về việc này thì tôi cũng không chắc, nhưng hiện tại tâm lý của tôi vẫn
còn ổn định và tỉnh táo. Tôi nghĩ tôi không điên, chỉ là có hơi sợ giao
tiếp với mọi người mà thôi.

Với lại, tôi cũng đã có một người bạn đang đợi ở nhà rồi.

Sau một thời gian khá dài ngồi trong lớp, cố gắng tiếp thu bài giảng của
giáo viên thì cuối cùng tiếng chuông của giờ tan trường cũng đã vang
lên. Tôi xách cặp, đi ra khỏi lớp và tiến thẳng về nhà. Dù sao tôi cũng
không có bạn nên cứ cố nán lại thì cũng không có ích gì. Với lại nếu về
trễ thì chị ấy sẽ thấy cô đơn mất.

``Con về rồi.'' Tôi cất tiếng chào khi vừa bước vào nhà.

Mà nói thế thôi chứ, trong nhà thật sự cũng chẳng có ai khác cả. Bố mẹ
vốn cũng chẳng yêu thương gì tôi cho lắm, họ chỉ có gửi tiền về cho tôi
đi học để làm tròn trách nhiệm và để chắc chắn sẽ không bị ăn kiện mà
thôi.

Tôi thật sự không cảm thấy buồn bã mấy vì suy cho cùng thì bản thân cũng
chẳng có tình mẫu hay phụ tử gì đó. Tất cả những gì tôi cần ở hai người
đó chỉ có số tiền được cấp hàng tháng mà thôi.

Với lại, tôi cũng có một gia đình khác rồi.

``Em về rồi đây!'' Tôi đẩy mạnh cửa phòng của mình.

``Mừng em về nhà, Yuzuha.'' Người đang ngồi ở bên trong vẫy tay chào
tôi.

Đây là chị Kiso, một người bạn đã chơi với tôi từ rất lâu trước kia.
Thật sự thì tôi cũng không nhớ rõ cả hai đã làm quen như thế nào. Chỉ
nhớ rằng chị Kiso đã chơi với tôi hồi mười năm trước và cả hai đã ở bên
nhau đến tận bây giờ.

Tôi và chị ấy đã luôn sống chung với nhau trong một căn phòng. Thỉnh
thoảng bố mẹ sẽ quay về đây để ngủ hay trú mưa, lúc đó tôi sẽ phải cố
gắng xoay sở làm sao để mà không cho họ phát hiện ra chị Kiso và đuổi
chị ấy đi mất.

``Hôm nay ở trường thế nào?'' Chị Kiso lên tiếng hỏi khi tôi đang thay
đồ, nhưng sau đó thì chị ấy cũng tự trả lời, ``Mà hỏi thế thôi chứ, chắc
lại chả thèm nói chuyện với ai phải không?''

``Chị kỳ quá đấy nhé, em cũng muốn nói chuyện với mọi người lắm mà.''
Tôi quay lại và phản bác lời nói của chị ấy. ``Chỉ là em không đủ can
đảm thôi.''

``Cũng đúng nhỉ?'' Chị Kiso nhìn tôi với một ánh mắt có phần đồng cảm,
sau đó thì mỉm cười. ``Mà, nếu không ai muốn nói chuyện với em thì cứ về
nhà với chị thôi. Chị sẽ lắng nghe em mà.''

``Chị Kiso...'' Tôi rưng rưng nước mắt và nhào đến định ôm lấy chị ấy.

``Này! Em quên lời chị nói rồi đấy hả?'' Chị ấy đứng dậy và né cái ôm,
khiến tôi đang lao đến va đầu vào cái tường trước mặt.

``Đau quá...'' Tôi xoa xoa cái đầu của mình. ``Chị là đồ xấu xa...''

``Do em cả thôi.'' Chị Kiso thở dài. ``Chẳng phải chị đã bảo rằng tuyệt
đối không được chạm vào chị rồi hay sao?''

``Thì đúng là như vậy, nhưng mà chị chả thèm nói cho em biết lý do gì
cả.'' Tôi bĩu môi phàn nàn.

Tuy là hai đứa tôi đã ở bên nhau rất lâu rồi nhưng tôi chưa bao giờ được
chạm vào người chị Kiso, và chị ấy cũng nghiêm cấm tôi làm việc đó. Chị
Kiso luôn rất thoải mái và chiều chuộng tôi, tôi muốn làm gì cũng được
nhưng chỉ có chuyện này là bị chị ấy cấm tiệt hoàn toàn.

Cái quy định bất thành văn này giữa cả hai đã tồn tại từ rất rất là lâu
trước kia rồi, nhưng đến giờ tôi vẫn không được biết lý do cho việc này.
Bộ ngay khi tôi chạm vào thì chị Kiso sẽ lăn đùng ra chết hay sao nhỉ?

``Khi nào cần thiết thì chị sẽ cho em biết, còn bây giờ thì vẫn chưa
được.''

``Tại sao chứ?'' Tôi ngẩng đầu lên hỏi.

``...Vì như thế sẽ tốt cho em hơn.''

Là sao cơ? Thỉnh thoảng chị ấy cứ cố làm ra cái vẻ bí bí ẩn ẩn để làm gì
vậy không biết nữa.

``Mà cũng có sao đâu.'' Như nhìn ra sự bất mãn hiện rõ trên gương mặt
tôi, chị Kiso lên tiếng an ủi. ``Dù không chạm vào, chị vẫn sẽ ở bên em
mà.''

``Hầy.''

Nếu chị ấy mà nói thế thì thôi vậy. Chị Kiso là gia đình duy nhất của
tôi, nếu mà cứ ép buộc thì chị ấy có khi sẽ đi mất.

Dù sao thì, tôi cũng chỉ còn mỗi chị ấy mà thôi.

*

``Giờ nhớ lại, đúng là mình đã từng có một thời kỳ như thế nhỉ?'' Tôi
ngồi trong phòng làm việc, bâng quơ nhớ lại chuyện cũ giữa giờ giải lao.

Chà, hồi trước đúng là tôi tự kỷ thật đấy ha? Tự tưởng tượng ra một
người bạn xong tự chơi chung luôn, đúng là hiện tại chả thể nghĩ được
rằng mình của hồi xưa lại như thế.

``Chị làm xong rồi thì ra ăn nhé.'' Cánh cửa đột nhiên mở ra, sau đó là
gương mặt thân thương của vợ tôi là Mirai. ``Em nấu cơm xong rồi đấy.''

``Ừm, chị ra ngay!'' Tôi đáp lại em ấy và bắt tay vào viết những chữ
cuối cùng.

Chà, nhưng dù khó tin cỡ nào đi nữa thì đó cũng là một khoảng thời gian
khá đẹp. Và cũng không thể phủ nhận rằng tôi của ngày xưa đã từng thiếu
thốn tình cảm nhiều đến như vậy.

Nhưng mà bây giờ thì khác rồi.

``Em đã có một cô vợ hết sức đáng yêu, có nghề nghiệp ổn định, có bạn
bè, không còn u ám và sợ xã hội như nữa. Thế nên...'' Tôi đóng máy tính
lại và đứng lên. ``Em đã ổn rồi.''

\emph{``Đừng quên chị nhé.''}

``Hơ?'' Tôi quay mặt ra sau lưng, nơi vừa phát ra giọng nói ấy.

Và sau đó, tôi mỉm cười với người từng là gia đình duy nhất của mình.

``Em tuyệt đối không bao giờ quên chị đâu, chị Kiso.''
}